\thispagestyle{empty}
\null
\vspace*{\fill}
\begin{center}
    {\LARGE \textbf{LAMPIRAN}}
\end{center}
\addcontentsline{toc}{chapter}{Lampiran} 
\vspace*{\fill}
\clearpage
\pagestyle{plain}

\chapter{Komentar Penguji dan Respon terhadap Komentar}

\section*{1. Komentar Pak Helmy}
\textbf{Komentar:} \\
Rumusan masalah pertama dibuat lebih lugas dan langsung menyebutkan permasalahan utama yang berkaitan dengan Reinforcement Learning.

\textbf{Respon:} \\
Rumusan masalah telah direvisi dengan menggunakan kalimat yang lebih lugas dan berfokus langsung pada penerapan \textit{Reinforcement Learning} dalam sistem \textit{automated web penetration testing}. Perubahan dilakukan pada bagian \textit{Rumusan Masalah} di Bab I. Kalimat sebelumnya yang terlalu panjang dan umum telah diganti menjadi pembahasan yang lebih spesifik mengenai penerapan RL untuk menentukan strategi pengujian kerentanan secara adaptif.

\vspace{0.5cm}

\section*{2. Komentar Pak Helmy}
\textbf{Komentar:} \\
Pada metode penelitian, skenario pengujian dibuat lebih detail.

\textbf{Respon:} \\
Bagian metode penelitian telah diperbaiki dengan menambahkan penjelasan yang lebih rinci mengenai skenario pengujian. Revisi mencakup penjelasan environment yang digunakan, definisi \textit{state}, \textit{action}, \textit{reward}, parameter \textit{Reinforcement Learning}, serta metode evaluasi sistem. Perubahan dapat dilihat pada Bab III bagian \textit{Metode Penelitian} dan subbab \textit{Skenario Pengujian Sistem}.

\vspace{0.5cm}

\section*{3. Komentar Pak Helmy}
\textbf{Komentar:} \\
Penelitian terkait nomor 8 kosong, dilengkapi atau dihapus apabila memang tidak ada.

\textbf{Respon:} \\
Bagian penelitian terkait yang kosong telah ada, dikarenakan kesalahan teknis yang ada sehingga bagian tersebut bisa hilang. Perubahan dapat dilihat pada Bab II bagian \textit{Penelitian Terdahulu}.

\vspace{0.5cm}

\section*{4. Komentar Pak Helmy}
\textbf{Komentar:} \\
Pembahasan yang dikerjakan banyak, tetapi isi pembahasannya masih terlalu sedikit.

\textbf{Respon:} \\
Isi pembahasan telah diperluas dengan menambahkan penjelasan teknis yang lebih mendalam terkait proses \textit{Reinforcement Learning}, mekanisme pengujian otomatis, serta integrasi sistem dengan OWASP Juice Shop. Selain itu, beberapa subbab yang terlalu singkat telah digabung agar pembahasan menjadi lebih fokus dan substantif. Perubahan dapat dilihat pada Bab II dan Bab III.

\vspace{0.5cm}

\section*{5. Komentar Pak Helmy}
\textbf{Komentar:} \\
Penggunaan tanda poin jangan berbentuk bullet, diganti menggunakan angka atau huruf.

\textbf{Respon:} \\
Seluruh format \textit{bullet point} pada dokumen telah diubah menjadi format enumerasi menggunakan angka dan huruf sesuai kaidah penulisan ilmiah. Perubahan dapat dilihat pada seluruh bagian proposal yang menggunakan daftar poin, khususnya pada Bab II dan Bab III.

\vspace{0.5cm}

\section*{6. Komentar Pak Helmy}
\textbf{Komentar:} \\
Caption tabel seharusnya berada di atas tabel.

\textbf{Respon:} \\
Seluruh \textit{caption} tabel telah dipindahkan ke bagian atas tabel sesuai format penulisan karya ilmiah yang benar. Perubahan dapat dilihat pada seluruh tabel dalam proposal, khususnya pada tabel penelitian terdahulu dan tabel parameter pengujian.

\vspace{0.5cm}

\section*{7. Komentar Pak Helmy}
\textbf{Komentar:} \\
Pemberian nilai reward ditentukan berdasarkan apa dan perlu diberikan acuannya.

\textbf{Respon:} \\
Penjelasan mengenai dasar penentuan nilai reward telah ditambahkan dengan mengacu pada konsep \textit{reward shaping} dalam \textit{Reinforcement Learning}. Selain itu, ditambahkan tabel reward beserta alasan pemberian nilai untuk setiap kondisi pengujian. Perubahan dapat dilihat pada Bab III bagian \textit{Perancangan Reward}.

\vspace{0.5cm}

\section*{8. Komentar Bu Qori}
\textbf{Komentar:} \\
Acuan penentuan nilai reward perlu dijelaskan lebih detail.

\textbf{Respon:} \\
Penjelasan mengenai mekanisme reward telah diperjelas dengan menambahkan referensi pendukung serta tabel kondisi reward dan penalti yang digunakan dalam penelitian. Revisi ini dapat dilihat pada Bab III subbab \textit{Perancangan Reward dan Punishment}.

\vspace{0.5cm}

\section*{9. Komentar Bu Qori}
\textbf{Komentar:} \\
Alur metode penelitian masih terlalu sederhana.

\textbf{Respon:} \\
Diagram dan penjelasan alur metode penelitian telah diperbaiki menjadi lebih rinci dengan menambahkan tahapan studi literatur, perancangan environment, implementasi RL agent, pelatihan agen, hingga evaluasi sistem. Perubahan dapat dilihat pada Bab III bagian \textit{Alur Penelitian}.

\vspace{0.5cm}

\section*{10. Komentar Bu Qori}
\textbf{Komentar:} \\
Bagian pembangunan prototipe sebaiknya dijadikan subbab di dalam perancangan sistem.

\textbf{Respon:} \\
Struktur penulisan telah direvisi dengan memindahkan bagian pembangunan prototipe ke dalam subbab perancangan sistem agar alur pembahasan menjadi lebih terstruktur. Perubahan dapat dilihat pada Bab III bagian \textit{Perancangan Sistem}.

\vspace{0.5cm}

